\documentclass[a4paper]{arrowhead}

\usepackage[yyyymmdd]{datetime}
\usepackage{etoolbox}
\usepackage[utf8]{inputenc}
\usepackage[T1]{fontenc}
\usepackage{multirow}
\usepackage{longtable}
\usepackage{tikz}
\usetikzlibrary{positioning,shapes.multipart,arrows.meta}

\renewcommand{\dateseparator}{-}
\setlength{\emergencystretch}{3em}
\lstset{language=HTTP,breakatwhitespace=false,columns=fullflexible}

%% Special references
\newcommand{\mref}[2][]{{\textcolor{ArrowheadPurple}{\hyperref[sec:model:#2]{\ifstrempty{#1}{#2}{#1}}}}}
\newcommand{\oref}[1]{{\textcolor{ArrowheadBlue}{\hyperref[sec:operations:#1]{#1}}}}
\newcommand{\pdef}[1]{{\textcolor{ArrowheadGrey}{#1\label{sec:model:primitives:#1}\label{sec:model:primitives:#1s}\label{sec:model:primitives:#1es}}}}
\newcommand{\pref}[1]{{\textcolor{ArrowheadGrey}{\hyperref[sec:model:primitives:#1]{#1}}}}

% as template_sd.tex; the input and output arguments are passed already formatted (\mref or \pref)
\newrobustcmd\osubsection[3]{
  \addtocounter{subsection}{1}
  \addcontentsline{toc}{subsection}{\protect\numberline{\thesubsection}operation \textcolor{ArrowheadBlue}{#1}}
  \subsection*{
    \thesubsection\quad
    operation
    \textcolor{ArrowheadBlue}{#1}
    (\notblank{#2}{#2}{})
    \notblank{#3}{: #3}{}
  }
  \label{sec:operations:#1}
}
\newrobustcmd\msubsection[2]{
  \addtocounter{subsection}{1}
  \addcontentsline{toc}{subsection}{\protect\numberline{\thesubsection}#1 \textcolor{ArrowheadPurple}{#2}}
  \subsection*{\thesubsection\quad#1 \textcolor{ArrowheadPurple}{#2}}
  \label{sec:model:#2} \label{sec:model:#2s} \label{sec:model:#2es}
}
%%

\usepackage{xurl}
\begin{document}
\sloppy

%% Arrowhead Document Properties
\ArrowheadTitle{dataHistorian}
\ArrowheadServiceID{dataHistorian}
\ArrowheadType{Service Description}
\ArrowheadTypeShort{SD}
\ArrowheadVersion{5.3.0}
\ArrowheadDate{2026-11-09}
\ArrowheadAuthor{Jerker Delsing}
\ArrowheadStatus{REVIEW DRAFT}
\ArrowheadContact{jerker.delsing@ltu.se}
\ArrowheadFooter{\href{www.arrowhead.eu}{www.arrowhead.eu}}
\ArrowheadSetup
%%

%% Front Page
\begin{center}
  \vspace*{1cm}
  \huge{\arrowtitle}

  \vspace*{0.2cm}
  \LARGE{\arrowtype}
  \vspace*{1cm}

  \Large{Service ID: \textit{"\arrowid"}}

  \vspace*{\fill}

  % Front Page Image

  \vspace*{1cm}
  \vspace*{\fill}

  % Front Page Abstract
  \begin{abstract}
    This document describes the \texttt{dataHistorian} service, version 1.0.0, produced by the DataManager support system of Arrowhead (documentation edition 5.3.0).
  \end{abstract}

  \vspace*{1cm}
\end{center}

\newpage
%%

%% Table of Contents
\tableofcontents
\newpage
%%

\section{Documentation edition and implementation scope}
This is the 5.3.0 documentation edition, dated 2026-11-09. It carries forward the system and service contracts from the retained 5.2.2 documentation and identifies the current primary implementation sources. The former \path|documentation/v5.2.2| folder is retained unchanged.

Document version 5.3.0 is distinct from executable, dependency, registered system/service and wire-contract versions. Values still defined by the source (including 5.2.2 artefact versions, 1.0.0 service versions and \path|SysMonitor-AMQP-IDD-5.2.2| payload identifiers) are preserved. This documentation update does not change runtime behaviour or establish a new implementation release. Historical verification statements describe the tests recorded in the inherited baseline; this revision validates document compilation and references.

\section{Overview}
\label{sec:overview}

This document describes the \texttt{dataHistorian} service, which provides persistent storage and query of SenML time series per system and service.
It is produced by the DataManager system \cite{dm-sysd} in service version 1.0.0.

The rest of this document is organized as follows.
In Section \ref{sec:operations}, we describe the abstract message operations provided by the service.
In Section \ref{sec:model}, we end the document by presenting the data types used by the mentioned operations.
The concrete interface is specified in the HTTP/TLS/JSON IDD \cite{dm-dataHistorian-http-idd}.

\subsection{Significant Prior Art}

Successor of the Arrowhead 4.6 DataManager historian service.

\subsection{How This Service Is Meant to Be Used}

Producers append records with \texttt{put-\allowbreak{}data}; consumers query with \texttt{get-\allowbreak{}data}, optionally limited by \texttt{from}, \texttt{to}, \texttt{count}, \texttt{signals} and \texttt{signalCount}.

\subsection{Important Delimitations and Dependencies}
\label{sec:delimitations}

Data are stored in DataManager's own database without retention policy. With \texttt{service.\allowbreak{}authorization.\allowbreak{}enabled=\allowbreak{}true} every call is checked with Consumer\allowbreak{}Authorization.

\newpage

\section{Service Interface}
\label{sec:operations}

This section describes the interfaces to the \texttt{dataHistorian} service.
In particular, each subsection names an abstract operation, an input type and an output type, in that order.
The input type is named inside parentheses, while the output type is preceded by a colon.
Input and output types are only denoted when accepted or returned, respectively, by the interface in question.

All abstract data types named in this section are defined in Section \ref{sec:model}.

\begin{figure}[ht!]
  \centering
  \begin{tikzpicture}[font=\small]
    \node[draw=ArrowheadBlue, very thick, rectangle split, rectangle split parts=2, align=center,
          rectangle split part fill={ArrowheadBlue!12,white}, text width=9cm, inner sep=6pt] (b)
      {\guillemotleft{}Service\guillemotright{}\\\textbf{dataHistorian}
        \nodepart{second} \texttt{list-\allowbreak{}systems}\\ \texttt{list-\allowbreak{}services}\\ \texttt{get-\allowbreak{}data}\\ \texttt{put-\allowbreak{}data}};
  \end{tikzpicture}
  \caption{SysML-style block of the \texttt{dataHistorian} service and its operations.}
  \label{fig:SD}
\end{figure}

The following interface operations are available.

\osubsection{list-systems}{}{\mref{SystemListResponse}}

Returns the systems with stored data.

\osubsection{list-services}{\pref{Name}}{\mref{ServiceListResponse}}

Returns the services of a system with stored data.

\osubsection{get-data}{\pref{Name}}{\pref{List}$<$\allowbreak{}\mref{SenMLRecord}$>$}

Returns stored SenML records of the system and service, newest first, limited by the query parameters.

\osubsection{put-data}{\pref{List}$<$\allowbreak{}\mref{SenMLRecord}$>$}{}

Appends SenML records for the system and service.

\newpage

\section{Information Model}
\label{sec:model}

Here, all data objects that can be part of the \texttt{dataHistorian} service are listed in alphabetic order.
Note that each subsection, which describes one type of object, begins with the \textit{struct} keyword, which is used to denote a collection of named fields, each with its own data type.
As a complement to the explicitly defined types in this section, there is also a list of implicit primitive types in Section \ref{sec:model:primitives}, which are used to represent things like hashes and identifiers.

\msubsection{struct}{SenMLRecord}

A SenML record (RFC 8428); a message is a JSON array of records.


\noindent\begin{tabularx}{\textwidth}{| >{\raggedright\arraybackslash}p{3.9cm} | >{\raggedright\arraybackslash\small}p{5.3cm} | X |} \hline
\rowcolor{gray!33} Field & Type & Description \\ \hline
bn & \pref{String} & SenML base name. \\ \hline
bt & \pref{String} & SenML base time (seconds). \\ \hline
bu & \pref{String} & SenML base unit. \\ \hline
bv & \pref{String} & SenML base value. \\ \hline
bs & \pref{String} & SenML base sum. \\ \hline
bver & \pref{String} & SenML base version. \\ \hline
n & \pref{String} & SenML name. \\ \hline
u & \pref{String} & SenML unit. \\ \hline
v & \pref{String} & SenML numeric value. \\ \hline
vs & \pref{String} & SenML string value. \\ \hline
vb & \pref{Boolean} & SenML Boolean value. \\ \hline
vd & \pref{String} & SenML data value (Base64). \\ \hline
s & \pref{String} & SenML sum. \\ \hline
t & \pref{String} & SenML time (seconds, relative to the base time). \\ \hline
ut & \pref{String} & SenML update time (seconds). \\ \hline
\end{tabularx}



\msubsection{struct}{ServiceListResponse}

Service names with data of a system.


\noindent\begin{tabularx}{\textwidth}{| >{\raggedright\arraybackslash}p{3.9cm} | >{\raggedright\arraybackslash\small}p{5.3cm} | X |} \hline
\rowcolor{gray!33} Field & Type & Description \\ \hline
services & \pref{List}$<$\allowbreak{}\pref{String}$>$ & Service names with stored data. \\ \hline
\end{tabularx}



\msubsection{struct}{SystemListResponse}

List of systems.


\noindent\begin{tabularx}{\textwidth}{| >{\raggedright\arraybackslash}p{3.9cm} | >{\raggedright\arraybackslash\small}p{5.3cm} | X |} \hline
\rowcolor{gray!33} Field & Type & Description \\ \hline
systems & \pref{List}$<$\allowbreak{}\pref{String}$>$ & System names with stored data. \\ \hline
\end{tabularx}



\subsection{Primitives}
\label{sec:model:primitives}

Types and structures mentioned throughout this document that are assumed to be available to implementations of this service.
The concrete interpretations of each of these types and structures must be provided by any IDD document claiming to implement this service.

\noindent\begin{tabularx}{\textwidth}{| p{4.6cm} | X |} \hline
\rowcolor{gray!33} Type & Description \\ \hline
\pdef{Boolean} & One out of \texttt{true} or \texttt{false}. \\ \hline
\pdef{List} & An \textit{array} of a known number of items, each having type A. \\ \hline
\pdef{Name} & A string identifier that is intended to be both human and machine-readable. \\ \hline
\pdef{String} & A chain of characters. \\ \hline
\end{tabularx}

\newpage

\bibliographystyle{IEEEtran}
\bibliography{bibliography}

\newpage

\section{Revision History}
\subsection{Amendments}

\noindent\begin{tabularx}{\textwidth}{| p{1cm} | p{3cm} | p{2cm} | X | p{4cm} |} \hline
\rowcolor{gray!33} No. & Date & Version & Subject of Amendments & Author \\ \hline

1 & 2026-10-22 & 5.2.2 & Document created from the Arrowhead LaTeX templates for the 5.2.2 Java Spring implementation, including the local extensions of this workspace. & Jerker Delsing \\ \hline
2 & 2026-10-30 & 5.2.2 & Checked against the primary Java source tree (Eclipse-Arrowhead-java): services, operations, paths and data models unchanged. & Claude Opus 5.5 (draft) \\ \hline
3 & 2026-11-09 & 5.3.0 & New documentation edition; primary-source references and Go/Rust implementation scope; 5.2.2 retained. & Documentation maintenance \\ \hline
\end{tabularx}

\subsection{Quality Assurance}

\noindent\begin{tabularx}{\textwidth}{| p{1cm} | p{3cm} | p{2cm} | X |} \hline
\rowcolor{gray!33} No. & Date & Version & Approved by \\ \hline

1 & & \arrowversion  &  \\ \hline

\end{tabularx}

\end{document}
